\section{Pairwise obstructions and certified forecasting}
\label{sec:theory}

This section separates three questions: whether a representation can support a
specified accuracy, whether a fitted forecaster attains that accuracy, and what
can be concluded when the representation is refined.  All population risks
refer to the same raw forecasting target.  No Markov assumption, invariant
distribution, or exact Koopman closure is imposed.

\subsection{A directly estimable pairwise obstruction}

Let $X$ be an available history, $Y\in\mathbb R^q$ a future target, and
$U=\Gamma(X)$ a fixed representation.  Let $d$ be a pseudometric on the
representation space and write
\[
 \mathcal G_\kappa(d)
 =\{g:\|g(u)-g(v)\|_2\leq\kappa d(u,v)\},\qquad
 R^*_\kappa(\Gamma,d)=\inf_{g\in\mathcal G_\kappa(d)}
             \mathbb E\|Y-g(U)\|_2^2.
\]
The metric, representation, and nonnegative budget $\kappa$ are fixed before
the certificate data are observed.  The infimum may be over a larger class
than the fitted decoder family.  A lower bound for this larger class is also
a lower bound for every decoder in the family.

For a conditioning variable $A$, construct $(U,Y)$ and $(U',Y')$ independently
conditional on $A$, with the same conditional law and with the original
marginal law for each copy.  The case of a constant $A$ is an ordinary
independent pair.  Define
\begin{equation}
 W_{\kappa,A}(\Gamma,d)
 =\frac12\mathbb E\left[
       \bigl(\|Y-Y'\|_2-\kappa d(U,U')\bigr)_+^2\right].
 \label{eq:pair-witness}
\end{equation}
The shared variable specifies a sampling construction; arbitrary nearest
neighbor matching is not covered by this definition.

\begin{theorem}[Pairwise prediction obstruction]
\label{thm:pair-obstruction}
Assume $\mathbb E\|Y\|_2^2<\infty$.  Then
\begin{equation}
 R^*_\kappa(\Gamma,d)\geq W_{\kappa,A}(\Gamma,d).
 \label{eq:pair-population}
\end{equation}
For Euclidean representation distance, assume also $\mathbb E\|U\|_2^2<\infty$ and put
\[
 \sigma_{Y\mid A}^2=\mathbb E\|Y-\mathbb E[Y\mid A]\|_2^2,
 \qquad
 \tau_{U\mid A}^2=\mathbb E\|U-\mathbb E[U\mid A]\|_2^2.
\]
Then the pairwise obstruction satisfies
\begin{equation}
 W_{\kappa,A}(\Gamma,\|\cdot\|_2)
 \geq (\sigma_{Y\mid A}-\kappa\tau_{U\mid A})_+^2.
 \label{eq:pair-dominance}
\end{equation}
Both comparisons use the same conditioning variable and evaluation law.
\end{theorem}

\begin{proof}
Fix an admissible $g$ with finite risk and set $r=Y-g(U)$ and
$r'=Y'-g(U')$.  The triangle inequality and the Lipschitz condition give
\[
 (\|Y-Y'\|_2-\kappa d(U,U'))_+\leq\|r-r'\|_2.
\]
Conditional independence gives
\[
 \mathbb E[\|r-r'\|_2^2\mid A]
 =2\{\mathbb E[\|r\|_2^2\mid A]
                    -\|\mathbb E[r\mid A]\|_2^2\}.
\]
Taking expectations proves $W_{\kappa,A}\leq R(g)$; taking the infimum
proves \eqref{eq:pair-population}.  Infinite-risk predictors satisfy the
inequality trivially.  To prove \eqref{eq:pair-dominance}, write
$a=\|Y-Y'\|_2$ and $b=\|U-U'\|_2$.  Pointwise,
$a\leq(a-\kappa b)_++\kappa b$.  The triangle inequality in $L^2$ therefore
implies
\[
 \|(a-\kappa b)_+\|_{L^2}
 \geq(\|a\|_{L^2}-\kappa\|b\|_{L^2})_+
 =\sqrt2(\sigma_{Y\mid A}-\kappa\tau_{U\mid A})_+.
\]
Square and divide by two.
\end{proof}

The new statistic estimates a squared positive-part distance discrepancy
directly.  It does not estimate two variances separately and then subtract
their square roots.  This distinction can matter when retained input
variation is small.  It is not a claim of a new minimax estimation rate:
bounded-mean concentration is used below, and taking a square root of a
risk bound can still have slow behavior near zero.  In particular,
\eqref{eq:pair-dominance} does not imply that pairing at a common time origin
dominates a certificate conditioned on a different reference class $c(X)$.

\begin{proposition}[Exactness on a two-atom law]
\label{prop:two-atom}
Let $(U,Y)$ equal $(u_0,y_0)$ with probability $p\in(0,1)$ and
$(u_1,y_1)$ otherwise.  Use Euclidean input distance and an unconditioned
independent pair.  Then
\begin{equation}
 R^*_\kappa
 =W_\kappa
 =p(1-p)\bigl(\|y_1-y_0\|_2
                 -\kappa\|u_1-u_0\|_2\bigr)_+^2.
 \label{eq:two-atom-optimum}
\end{equation}
The optimum is attainable by an affine map.  If the outputs are constrained
to a closed convex set containing $y_0,y_1$, projecting that map onto the
set preserves the optimum and the Lipschitz budget.
\end{proposition}

\begin{proof}
Let $q=1-p$, $a=\|y_1-y_0\|_2$, $b=\|u_1-u_0\|_2$, and
$\mu=py_0+qy_1$.  The two unequal ordered pairs have total probability
$2pq$, so \eqref{eq:pair-witness} gives the right side of
\eqref{eq:two-atom-optimum} as a lower bound.  If $a=0$, a constant predictor
attains zero.  Otherwise let $v=(y_1-y_0)/a$ and $s=\min\{a,\kappa b\}$.
Choose fitted values $f_0=\mu-qsv$ and $f_1=\mu+psv$.  Their separation is
$s$, and their risk is $pq(a-s)^2$.  If $b=0$, these values coincide and
a constant map attains this risk.  If $b>0$, the affine map
\[
 g(u)=f_0+
 \frac{\langle u-u_0,u_1-u_0\rangle}{b^2}(f_1-f_0)
\]
has operator norm $s/b\leq\kappa$ and the prescribed two fitted values.
Both fitted values lie on the segment between $y_0$ and $y_1$.  Euclidean
projection onto a closed convex output set fixes these values and is
nonexpansive.
\end{proof}

\subsection{Independent episodes, dependent observations within episodes}

Let $E_1,\ldots,E_n$ be independent, identically distributed complete
episodes, independent of all model construction and selection data.  For
horizon $h$, let $\mathcal T_h$ be a fixed nonempty set of available causal
origins common to all episodes and all compared candidates.  Within an
episode, observations may have arbitrary dependence and their distribution
may change with time.  Candidate $j$ has a representation $\Gamma_j$,
metric $d_j$, frozen forecaster $f_{jh}=g_{jh}\circ\Gamma_j$, and known
Lipschitz budget $\kappa_{jh}$ for its decoder.  Denote its episode loss by
\begin{equation}
 Z_{i,jh}=\frac1{|\mathcal T_h|}\sum_{t\in\mathcal T_h}
       \|Y_{i,t+h}-f_{jh}(X_{i,t})\|_2^2,
 \qquad R_{jh}=\mathbb E Z_{i,jh}.
 \label{eq:episode-risk}
\end{equation}
This is risk for a new episode and a uniformly sampled common origin.  It
is not a guarantee conditional on each particular history.  Different
episode weights or a different origin law define a different risk.

Assume the target and prediction lie in a known closed convex output set
of diameter $D_h$, and put $B_h=D_h^2$.  Prediction clipping alone is
insufficient if the target is unbounded.  Bounds inferred only from observed
training minima and maxima do not establish the assumed population range.
For $m=\lfloor n/2\rfloor$ disjoint episode pairs define
\begin{equation}
 Q_{r,jh}=\frac1{2|\mathcal T_h|}\sum_{t\in\mathcal T_h}
 \left(\|Y_{2r-1,t+h}-Y_{2r,t+h}\|_2
 -\kappa_{jh}d_j(\Gamma_j(X_{2r-1,t}),
                         \Gamma_j(X_{2r,t}))\right)_+^2.
 \label{eq:episode-pairs}
\end{equation}
The $Z_i$ are independent across episodes, and the $Q_r$ are independent
across disjoint pairs.  The terms within either average need not be
independent.  Their ranges are $[0,B_h]$ and $[0,B_h/2]$, respectively.
For the pair statistic, the conditioning variable in
Theorem~\ref{thm:pair-obstruction} is the common random origin $A=t$.

For observations $V_1,\ldots,V_N$ in an interval of length $b$, let
$\bar V$ be their mean and $s_V^2=(N-1)^{-1}\sum_i(V_i-\bar V)^2$.
For $N\geq2$ define
\begin{equation}
 \operatorname{rad}_{\rm EB}(V;b,\alpha)
 =\sqrt{\frac{2s_V^2\log(2/\alpha)}{N}}
       +\frac{7b\log(2/\alpha)}{3(N-1)}.
 \label{eq:eb-radius}
\end{equation}
This is the bounded-variable empirical Bernstein radius of
\citet{maurer2009empirical}; applying it to the reflected variables gives
the lower confidence bound used here.

\begin{theorem}[Simultaneous two-sided episode certificate]
\label{thm:episode-certificate}
Suppose $J$ candidates and $H$ horizons are frozen using independent
development data, $n\geq4$, and all fitted decoders satisfy their stated
Lipschitz budgets.  Set $\alpha=\delta/(2JH)$ and
\begin{align}
 L_{jh}&=\max\{0,\bar Q_{jh}
 -\operatorname{rad}_{\rm EB}(Q_{\cdot,jh};B_h/2,\alpha)\},
 \label{eq:pair-lcb}\\
 U_{jh}&=\min\{B_h,\bar Z_{jh}
 +\operatorname{rad}_{\rm EB}(Z_{\cdot,jh};B_h,\alpha)\}.
 \label{eq:risk-ucb}
\end{align}
With probability at least $1-\delta$, conditional on the development data,
all $j,h$ simultaneously satisfy
\begin{equation}
 L_{jh}\leq W_{jh}\leq R^*_{jh}\leq R_{jh}\leq U_{jh}.
 \label{eq:risk-sandwich}
\end{equation}
Here $R^*_{jh}$ is the minimum risk over all $\kappa_{jh}$-Lipschitz
decoders of $\Gamma_j$ under the same horizon, origin law, and raw target.
\end{theorem}

\begin{proof}
Condition on development data.  For each fixed $j,h$, apply the empirical
Bernstein inequality to the $n$ episode losses and its reflected version
to the $m$ disjoint-pair scores.  A union bound over the $2JH$ events gives
simultaneous validity with failure probability at most $\delta$.
Theorem~\ref{thm:pair-obstruction}, applied at each common origin and
averaged over origins, gives $W_{jh}\leq R^*_{jh}$.  The fitted decoder
belongs to the infimum's class, giving $R^*_{jh}\leq R_{jh}$.  Independence
between the upper and lower statistics is not required.
\end{proof}

The confidence parameter $\delta$ controls the random calibration sample.
It does not state that every future error is at most $\sqrt U$, or give
conditional coverage of a prediction interval.  The effective independent
sample sizes in \eqref{eq:eb-radius} are $n$ and $m$, not the number of
overlapping windows.  More origins can reduce the variability of an
episode average, but cannot be counted as additional independent episodes.
The representation, score normalization, dictionary sizes, budgets,
horizons, and any baseline included in the reported joint event must all
be fixed before calibration.  Training-dependent choices are allowed
because the guarantee conditions on training.

\begin{corollary}[Acceptance, exclusion, and inconclusive decisions]
\label{cor:decisions}
On the event in \eqref{eq:risk-sandwich}, $U_{jh}\leq\epsilon_h^2$
certifies that the fitted forecaster attains root-mean-square error at
most $\epsilon_h$.  The inequality $L_{jh}>\epsilon_h^2$ excludes that
tolerance for every decoder in the corresponding Lipschitz class.
When neither condition holds, the decision is inconclusive.  Among a
frozen finite family, selecting any candidate satisfying
$U_{jh}\leq\epsilon_h^2$ at every required horizon preserves these
guarantees.  In particular, a predeclared cost can be minimized over the
certified candidates, with abstention if the set is empty.
\end{corollary}

The selected candidate is cheapest among those certified by the data,
not necessarily among all candidates whose unknown true risks meet the
tolerance.  Also, $\sqrt{U_{jh}}-\sqrt{L_{jh}}$ upper-bounds the fitted
decoder's RMSE suboptimality within the stated class, but a large gap
does not diagnose optimization error: the lower witness and both
confidence corrections can be loose.

For transparent sample-size calculations, the empirical Bernstein radii
can instead be replaced by the prespecified Hoeffding radii
$b\sqrt{\log(1/\alpha)/(2N)}$.  If a fitted risk is at least $\gamma>0$
below a squared-error tolerance, $n\geq
2B_h^2\log(1/\alpha)/\gamma^2$ makes its Hoeffding upper bound pass
that tolerance with probability at least $1-\alpha$.  Similarly, if a
pair-witness mean exceeds the tolerance by $\gamma$, then
$m\geq B_h^2\log(1/\alpha)/(2\gamma^2)$ makes the corresponding lower
bound exclude it with probability at least $1-\alpha$.  These statements
follow by bounding the empirical mean's deviation in the direction
opposite to the confidence bound.  They are sufficient sample sizes,
not optimal rates, and use the witness margin rather than the unknown
oracle-risk margin for exclusion.  Taking the better of separately
computed Hoeffding and empirical Bernstein bounds requires a joint
confidence allocation or a single prespecified choice.

\subsection{Refinement that preserves the meaning of the norm budget}

Let $\Phi_1,\ldots,\Phi_s$ be simplex-valued geometric dictionaries, with
possibly different numbers of coordinates, and retain the whole stack
$\Gamma_s=(\Phi_1,\ldots,\Phi_s)$.  Refinement appends a block; it does
not silently refit or replace previous blocks.  Use the block maximum
metric and the sum of block Frobenius norms,
\begin{equation}
 d_s(x,x')=\max_{1\leq b\leq s}
                  \|\Phi_b(x)-\Phi_b(x')\|_2,
 \qquad
 \|W\|_{2,1}=\sum_{b=1}^s\|W_b\|_{\rm F}.
 \label{eq:block-geometry}
\end{equation}
The notation $d_s$ denotes the induced pseudometric on histories, or
equivalently the block maximum metric on the stacked representation.
The decoder class is
\begin{equation}
 \mathcal F_{s,\kappa}
 =\left\{x\mapsto\Pi_{\mathcal Y}\!
       \left(a+\sum_{b=1}^sW_b\Phi_b(x)\right):
                         \sum_b\|W_b\|_{\rm F}\leq\kappa\right\},
 \label{eq:block-decoder}
\end{equation}
where the intercept $a$ is unpenalized and $\Pi_{\mathcal Y}$ is
Euclidean projection onto the convex output set.  Group norm constraints
are established tools \citep{bach2008group}; their role here is to align retained geometric
blocks with the sensitivity budget used in the certificate.

\begin{theorem}[Nested and replication-invariant refinement]
\label{thm:replication-invariance}
Every decoder in $\mathcal F_{s,\kappa}$ is $\kappa$-Lipschitz with
respect to $d_s$.  Appending a block preserves
$\mathcal F_{s,\kappa}\subseteq\mathcal F_{s+1,\kappa}$, and
$d_{s+1}\geq d_s$ pointwise.  Consequently, the optimal risk in the
decoder class cannot increase, and the population pair witness cannot
increase at the same $\kappa$ and with the same pairing law.  Appending
an exact copy of an existing block leaves both the decoder class and
the pair witness unchanged.
\end{theorem}

\begin{proof}
For the unprojected decoder,
\[
 \left\|\sum_bW_b\bigl(\Phi_b(x)-\Phi_b(x')\bigr)\right\|_2
 \leq\sum_b\|W_b\|_{\rm F}
                  \|\Phi_b(x)-\Phi_b(x')\|_2
 \leq\kappa d_s(x,x').
\]
Projection is nonexpansive, so the same bound holds afterward.  Set a
new block's coefficient to zero to embed any old decoder in the refined
class.  The maximum distance increases when a new block is added;
therefore each positive-part pair score decreases or stays unchanged.
If the new block duplicates block $b$, its contribution can be merged
into $W_b$ and
$\|W_b+W_{s+1}\|_{\rm F}\leq\|W_b\|_{\rm F}+
\|W_{s+1}\|_{\rm F}$.  This proves reverse inclusion of the decoder
classes.  A duplicate does not change a maximum distance.
\end{proof}

The risk statement concerns population optima in nested feasible sets,
not the performance of separately fitted finite-sample estimators.
Nor does it imply that every newly learned block supplies useful
information.  The maximum metric remains bounded by $\sqrt2$ because
each block is simplex-valued; its range does not grow merely by adding
blocks.  If a fitted decoder's smaller actual norm is used for an
additional certificate, that changes the comparator class.  Monotonicity
across refinements is claimed only at a common norm budget.

This distinction is material for a usual Frobenius constraint on the
concatenated coefficient matrix.  Consider a two-point input with
$\Phi(x_0)=e_1$, $\Phi(x_1)=e_2$, equally likely targets $-1$ and $1$,
and budget one.  Its best affine MSE is
$(2-\sqrt2)^2/4$.  Duplicating the block permits coefficients
$(-1/2,1/2,-1/2,1/2)$ of Frobenius norm one and zero error.  No information
was added: duplication only relaxed the effective norm constraint.
The block-sum construction prevents this effect.  It does not prevent
a genuinely different block from changing the useful geometry.

\begin{theorem}[A necessary amount of geometric change]
\label{thm:necessary-expansion}
Let $d_0,d_1$ be two induced representation pseudometrics on the same
histories, and use the same target, conditionally independent pairing
law, and positive sensitivity budget $\kappa$.  Write
$\Delta=(d_1(X,X')-d_0(X,X'))_+$ and let $W_0,W_1$ be their pair
witnesses.  Then
\begin{equation}
 \sqrt{W_1}\geq
 \left(\sqrt{W_0}-\frac{\kappa}{\sqrt2}
                       \|\Delta\|_{L^2}\right)_+.
 \label{eq:metric-change}
\end{equation}
In particular, if any $\kappa$-Lipschitz decoder of the new
representation attains RMSE at most $\epsilon$, then
\begin{equation}
 \|\Delta\|_{L^2}\geq
       \frac{\sqrt2}{\kappa}(\sqrt{W_0}-\epsilon)_+.
 \label{eq:required-change}
\end{equation}
On any event $L_0\leq W_0$, replacing $W_0$ by $L_0$ in the right side
of \eqref{eq:required-change} is valid simultaneously for every new
representation.  The new representations need not have been fixed
when $L_0$ was computed.
\end{theorem}

\begin{proof}
Let $a=\|Y-Y'\|_2$ and $v_i=(a-\kappa d_i)_+$.  The hinge function is
nonexpansive, and pointwise $v_0\leq v_1+\kappa\Delta$.
Consequently $\|v_1\|_{L^2}\geq
(\|v_0\|_{L^2}-\kappa\|\Delta\|_{L^2})_+$, which proves
\eqref{eq:metric-change}.  A feasible new decoder with RMSE at most
$\epsilon$ implies $W_1\leq\epsilon^2$ by
Theorem~\ref{thm:pair-obstruction}; rearrange.  The final statement
follows from the deterministic inequality and monotonicity in $W_0$.
\end{proof}

This is a necessary condition, not a sufficient refinement rule.  It
rules out repairs whose pairwise distance expansion is too small.  The
quantity $L_0$ must lower-bound the old \emph{pair witness}; an arbitrary
lower bound on old oracle risk cannot be substituted.  If sensitivity
budgets are also changed, the corresponding necessary condition is
\[
 \| (\kappa_1d_1-\kappa_0d_0)_+\|_{L^2}
 \geq\sqrt2(\sqrt{W_{\kappa_0,0}}-\epsilon)_+.
\]
This form makes explicit that increasing sensitivity can compensate for
changing the geometry.

\subsection{When an added block can improve prediction}

The preceding obstruction only identifies a necessary geometric change.
The following elementary quadratic-risk calculation identifies a
sufficient population condition for a particular residual direction.
It is a projection calculation, not a claim that a positive obstruction
guarantees a useful split.

\begin{proposition}[Residual correlation and budget slack]
\label{prop:residual-gain}
Consider the unprojected affine decoder
\[
 f_0(x)=a+\sum_{b=1}^sW_b\Phi_b(x),
\]
with $\sum_b\|W_b\|_{\rm F}<\kappa$, and let
$\psi(X)\in\mathbb R^p$ be a new square-integrable feature block.
Write
\[
 r=Y-f_0(X),\quad C=\mathbb E[r\psi(X)^\top],\quad
 G=\mathbb E[\psi(X)\psi(X)^\top],\quad
 s_0=\kappa-\sum_b\|W_b\|_{\rm F}.
\]
If $C\ne0$, define $a_0=\|C\|_{\rm F}^2$,
$b_0=\operatorname{tr}(CGC^\top)$, and
\[
 \eta=\min\{a_0/b_0,\ s_0/\sqrt{a_0}\}.
\]
Then $b_0>0$, the appended coefficient $V=\eta C$ is feasible, and
\begin{equation}
 R(f_0)-R(f_0+V\psi)
       =2\eta a_0-\eta^2b_0\geq\eta a_0>0.
 \label{eq:residual-gain}
\end{equation}
If $f_0(X)$ already lies in the convex output set almost surely,
projecting the new forecast onto that set retains this improvement.
\end{proposition}

\begin{proof}
Expanding the square gives
$R(f_0+V\psi)=R(f_0)-2\langle V,C\rangle_{
m F}
+\operatorname{tr}(VGV^\top)$.  For $V=\eta C$ this yields the equality
in \eqref{eq:residual-gain}.  The budget holds because
$\|V\|_{
m F}=\eta\sqrt{a_0}\leq s_0$.  If $b_0=0$, then
$C\psi=0$ almost surely.  For each row $c$ of $C$ this implies
$c\psi=0$ almost surely, whence
$\|c\|_2^2=\mathbb E[r_\ell\psi^\top]c^\top
=\mathbb E[r_\ell(c\psi)]=0$, contradicting $C\ne0$.
Since $\eta\leq a_0/b_0$, the gain is at least $\eta a_0$.
Projection onto a convex set containing $Y$ cannot increase squared
error; when $f_0$ is already in the set, its reference risk is unchanged.
\end{proof}

The proviso about the old output is essential.  Improvement of two
unprojected affine risks does not, by itself, imply improvement between
their separately clipped versions.  Nor does the proposition guarantee
strict improvement when the old decoder exhausts its budget.  For
actual clipped models, improvement is checked directly using paired
episode losses.  The moments in Proposition~\ref{prop:residual-gain} are
population moments: substituting estimates from training does not by
itself retain the claimed population gain.

For two forecasts fixed before an independent evaluation block, define
$D_i=Z_i^{\rm old}-Z_i^{\rm new}\in[-B_h,B_h]$.  With probability at
least $1-\alpha$,
\begin{equation}
 R^{\rm old}-R^{\rm new}\geq
 \bar D-\operatorname{rad}_{\rm EB}(D;2B_h,\alpha).
 \label{eq:paired-improvement}
\end{equation}
A positive right side certifies actual risk reduction for the two
fitted predictors.  A nonpositive value is inconclusive.  To certify
many fixed comparisons jointly, allocate confidence over the whole
comparison family.  For genuinely adaptive sequential refinement,
each new candidate may depend on the past, but its evaluation block
must be fresh and independent conditional on that past.  Assigning
stage failure probabilities $\alpha_s$ with
$\sum_s\alpha_s\leq\delta$ and accepting a change only when its
paired lower bound is positive then makes every accepted improvement
valid on a common event of probability at least $1-\delta$.  This
follows by conditional validity at each stage and a union bound; it
does not authorize repeatedly tuning on one fixed calibration set.

Finally, a positive obstruction need not be repairable.  If $Y$ is an
independent, symmetric random sign of amplitude $a$, then every
representation measurable from $X$ has Bayes risk $a^2$.  A constant
representation already has pair witness $a^2$; no refinement of the
observed history can remove this independent future noise.  Conversely,
a zero pair witness is not proof of a useful predictor.  Matched-origin
pairing, for example, sees no variation between identical deterministic
episodes even if an inadequate representation loses variation between
origins.  These examples motivate reporting obstruction, fitted risk,
and inconclusive outcomes separately.
